%% toc: P15 164 164.2
% Proofs for the modern reconstructions accompanying Milton's history of induction before Hume.
\begin{annexitem}{annex:p15-w02b-p15-026-conversion}{thm:p15-w02b-p15-026-conversion}{The conversion step in an extensional reconstruction}
\begin{proof}
(1) Let $x\in B$. By the convertibility premise $B\subseteq C$, $x\in C$.
By the premise $C\subseteq A$, $x\in A$. Since $x$ was arbitrary,
$B\subseteq A$. The premise $C\subseteq B$ is not used in this direction;
together with $B\subseteq C$ it makes $B=C$, which is the sense of
``convertible''.

(2) Take the domain $\{1,2\}$ and put $A=C=\{1\}$, $B=\{1,2\}$. Then
$C\subseteq A$ and $C\subseteq B$, but $2\in B\setminus A$, so
$B\not\subseteq A$. Here $B\not\subseteq C$: the convertibility premise
fails, and with it the inference. Nothing in the countercase depends on
whether the members of $C$ are read as individuals or as species.
\end{proof}
\end{annexitem}

\begin{annexitem}{annex:p15-w02b-p15-026-enumeration}{thm:p15-w02b-p15-026-enumeration}{Complete, covering and partial enumeration}
\begin{proof}
(1) Let $s\in S$. Since $S\subseteq O$, $s\in S\cap O$, and the evidence
$S\cap O\subseteq P$ gives $s\in P$. Hence $S\subseteq P$.

(2) Let $s\in S=K_1\cup\dots\cup K_m$. Then $s\in K_i$ for some $i$, and
$K_i\subseteq P$ gives $s\in P$; so $S\subseteq P$. For the second claim,
suppose $x\in S$. Then $x\in K_i$ for some $i$, contrary to the assumption
that $x$ is in none of the classes. Hence $x\notin S$. Neither claim needs
the classes to be disjoint or to be the finest division available; only
their union matters.

(3) Since $S\not\subseteq O$, choose $s_0\in S\setminus O$. Define
\[
P_1=(P\cap O)\cup(D\setminus O),\qquad P_2=P\cap O .
\]
Both are subsets of $D$, and $P_1\cap O=P\cap O=P_2\cap O$, because
$(D\setminus O)\cap O=\varnothing$. For $S\subseteq P_1$, take $s\in S$.
If $s\in O$, then $s\in S\cap O\subseteq P$, so $s\in P\cap O\subseteq
P_1$. If $s\notin O$, then $s\in D\setminus O\subseteq P_1$. For
$S\not\subseteq P_2$, note that $s_0\in S$ but $s_0\notin O$, so
$s_0\notin P\cap O=P_2$. Both $P_1$ and $P_2$ agree with $P$ on every
examined individual, so the evidence, which constrains $P$ only on $O$, is
satisfied by both, and it is compatible with the universal and with its
failure.
\end{proof}
\end{annexitem}

\begin{annexitem}{annex:p15-w02b-p15-026-ockham}{thm:p15-w02b-p15-026-ockham}{Ockham's rules}
\begin{proof}
(1) If ``every $S$ is $P$'' is true on the import reading, then
$S\neq\varnothing$ and $S\subseteq P$. Choose $s\in S$; then $s\in P$, so
the singular ``$s$ is $P$'' is true. On the reading $S\subseteq P$ alone,
the case $S=\varnothing$ makes the universal true while there is no
singular at all, so no true singular exists. If $S\neq\varnothing$, the
argument just given applies unchanged, so the failure occurs exactly when
$S=\varnothing$.

(2) Suppose every singular ``$s$ is $P$'' with $s\in S$ is true, that is,
$S\subseteq P$. If $S\neq\varnothing$, both conditions of the import
reading hold and the universal is true. If $S=\varnothing$, the hypothesis
holds vacuously, but the universal affirmative is false on the import
reading because its subject is empty.

(3) If ``no $S$ is $P$'' is false, then $S\cap P\neq\varnothing$. Choose
$s\in S\cap P$. The singular ``$s$ is not $P$'' is false, since $s\in P$.

(4) Let $p$ be a contingent proposition, which exists by hypothesis. By
the definition of contingency, $p$ can be true, so the singular ``$p$, a
contingent proposition, is true'' is possible. The same holds for every
contingent proposition, so every singular is possible. Suppose, for
contradiction, that ``every contingent proposition is true'' is possible.
Fix an alternative in which it holds. By hypothesis the contradictory
$\bar p$ of $p$ is contingent, and the class of contingent propositions is
the same in that alternative. So both $p$ and $\bar p$ fall under the
universal there and both are true there. This contradicts the hypothesis
that no proposition is true together with its contradictory. Hence the
universal is not possible, although each of its singulars is.
\end{proof}
\end{annexitem}

\begin{annexitem}{annex:p15-w02b-p15-026-exclusion}{thm:p15-w02b-p15-026-exclusion}{Exclusion under limited variety}
\begin{proof}
(0) Suppose $n$ agrees with $g$ on $I$. If (E1) applied at some $i\in I$,
then $\delta_g(i)>0$ and $\delta_n(i)=0$, contradicting (A1) at $i$. If
(E2) applied at some $i\in I$, then $\delta_g(i)=0$ and $\delta_n(i)>0$,
again contradicting (A1) at $i$. If (E3) applied at some $i,j\in I$, then
$\delta_g(j)<\delta_g(i)$ and $\delta_n(j)>\delta_n(i)$, contradicting (A2)
for $i,j$. So $n$ survives.

Conversely, suppose $n$ does not agree with $g$ on $I$. If (A1) fails at
some $i\in I$, then exactly one of $\delta_n(i)>0$ and $\delta_g(i)>0$
holds. Degrees are non-negative, so either $\delta_g(i)>0$ and
$\delta_n(i)=0$, which is (E1), or $\delta_g(i)=0$ and $\delta_n(i)>0$,
which is (E2). If (A2) fails for some $i,j\in I$, then
$\delta_g(j)<\delta_g(i)$ and not $\delta_n(j)\leq\delta_n(i)$, that is,
$\delta_n(j)>\delta_n(i)$, which is (E3). In either case $n$ is excluded.

(1) Let $f\in\mathcal N$ be a form of $g$. Then $f$ agrees with $g$ on
$\Omega$. Conditions (A1) and (A2) for $J=\Omega$ quantify over all
members of $\Omega$, so they hold in particular for all members of
$I\subseteq\Omega$. Hence $f$ agrees with $g$ on $I$, and $f$ survives by
(0).

(2) By (1), $f$ is among the survivors. If exactly one candidate
survives, that candidate is $f$.

(3) By (0), each of the two candidates survives. So at least two
candidates survive, and the surviving set, which contains both, does not
distinguish one of them from the other.

(4) Take $\Omega=\{i,i'\}$ with $i\neq i'$ and $I=\{i\}$. Put
$\delta_g(i)=\delta_g(i')=1$. Let $\mathcal N=\{n\}$ with $\delta_n(i)=1$
and $\delta_n(i')=0$.

The candidate $n$ agrees with $g$ on $I$. Condition (A1) at $i$ holds
because $\delta_n(i)=1>0$ and $\delta_g(i)=1>0$. Condition (A2) on $I$
concerns only the pair $(i,i)$, and its antecedent
$\delta_g(i)<\delta_g(i)$ is false, so it holds. By (0), $n$ survives.
Since $\mathcal N=\{n\}$, $n$ is the unique survivor.

The candidate $n$ is not a form of $g$. At $i'\in\Omega$ we have
$\delta_g(i')=1>0$ and $\delta_n(i')=0$, so (A1) fails on $\Omega$. As $n$
is the only candidate, no member of $\mathcal N$ is a form of $g$, and the
form hypothesis fails. The unique survivor is therefore not a form of $g$,
so the conclusion of part~(2) fails in this case.

For the last claim, let $I=\Omega$ and suppose no member of $\mathcal N$
is a form of $g$. Take any $n\in\mathcal N$. It does not agree with $g$
on $\Omega$, and $\Omega=I$, so it does not agree with $g$ on $I$. By (0),
$n$ is excluded. Hence a candidate that is not a form of $g$ can survive
only when $\Omega\setminus I$ is non-empty, as it is in the example, where
$\Omega\setminus I=\{i'\}$.
\end{proof}
\end{annexitem}

\begin{annexitem}{annex:p15-w02b-p15-026-pascal}{thm:p15-w02b-p15-026-pascal}{The table rule for every power}
\begin{proof}
(1) By definition $(a+e)^n=\sum_{k=0}^{n}c(n,k)\,a^{n-k}e^k$, and the
convention $c(n,k)=0$ outside $0\leq k\leq n$ lets the sum run over all
integers $k$. Multiplying by $a+e$ and using commutativity and
distributivity,
\[
(a+e)^{n+1}=\sum_{k}c(n,k)\,a^{n+1-k}e^{k}+\sum_{k}c(n,k)\,a^{n-k}e^{k+1}.
\]
Re-indexing the second sum by $k\mapsto k-1$ gives
$\sum_k c(n,k-1)\,a^{n+1-k}e^{k}$. Hence the coefficient of
$a^{n+1-k}e^k$ in $(a+e)^{n+1}$ is $c(n,k)+c(n,k-1)$. Monomials
$a^{n+1-k}e^k$ with different $k$ are distinct, so this coefficient is
$c(n+1,k)$. For $k<0$ or $k>n+1$ both sides are $0$. This argument is
carried out for an arbitrary fixed $n$ and uses no hypothesis about other
values of $n$.

(2) By induction on $n$. For $n=0$, $(a+e)^0=1$, so $c(0,0)=1=t(0,0)$, and
both vanish for $k\neq 0$. Suppose $t(n,k)=c(n,k)$ for all $k$. Then for
every $k$,
\[
t(n+1,k)=t(n,k-1)+t(n,k)=c(n,k-1)+c(n,k)=c(n+1,k),
\]
the last step by (1). Hence $t(n+1,k)=c(n+1,k)$ for all $k$, and the claim
holds for all $n\geq 0$. For $n\leq 5$ the table gives the rows printed
in Milton's chapter, and for $n=2,3,4$ these agree with the expansions
shown there.
\end{proof}
\end{annexitem}
